\chapter{Infrared Divergences and Inclusive QED}
\label{ch:volume-ii-infrared-inclusive-qed}
\ConventionsLorentzian


\section{Soft Radiation Data}

The scattering constructions developed earlier are simplest when all
asymptotic particles are massive.  Four-dimensional QED has a massless stable
photon, and a charged
particle may emit photons of arbitrarily small energy.  This chapter constructs
the scattering quantities adapted to that long-distance structure.

Throughout the chapter, spacetime is \(D=4\) Minkowski space with metric
\[
  \eta_{\mu\nu}=\operatorname{diag}(-1,+1,+1,+1).
\]
Consider a hard transition
\[
  \alpha\longrightarrow\beta
\]
between massive charged particles.  The hard external charged lines are
labelled by \(n\).  Their momenta are \(p_n\), their charges are \(g_n\), and
\[
  p_n^2=-m_n^2,\qquad p_n^0>0 .
\]
The sign
\[
  \eta_n=
  \begin{cases}
    +1,& n\in\beta\quad\hbox{outgoing},\\
    -1,& n\in\alpha\quad\hbox{incoming}
  \end{cases}
\]
records whether the line is outgoing or incoming.  Charge conservation for
the hard process is
\[
  \sum_n \eta_n g_n=0 .
\]
The charge symbols used in the QED chapters are related as follows.
\begin{center}
\begin{tabular}{p{0.24\textwidth}|p{0.12\textwidth}|p{0.54\textwidth}}
context & symbol & meaning \\
\hline
Chapter~\ref{ch:qed-external-states}
  & \(g\)
  & signed Abelian charge in \(D_\mu=\partial_\mu-\ii gA_\mu\);
    for an electron \(g=-e\), with \(e>0\). \\
Chapter~\ref{ch:qed-renormalization-form-factors}
  & \(e,\ e^B\)
  & positive magnitude of the electron charge and its bare counterpart; the
    electron covariant derivative is \(\partial_\mu+\ii eA_\mu\), and
    \(e=e^B Z_A^{1/2}\). \\
this chapter
  & \(g_n\)
  & signed physical charge of the \(n\)-th hard external line; an electron has
    \(g_n=-e\), a positron has \(g_n=+e\), and an external scalar of charge
    \(g\) in the scalar-QED example has \(g_n=g\).
\end{tabular}
\end{center}
Thus \(\eta_n\) records whether the charged line is incoming or outgoing,
whereas \(g_n\) records the physical sign of the charge carried by that line.

A soft photon has momentum
\[
  k^\mu=(\omega,\vec k),\qquad \omega=|\vec k|,
\]
and physical helicity \(h=\pm\).  The detector threshold is denoted by
\(E_T\): photons below this energy are included in the same observed hard
event.  An auxiliary infrared cutoff \(\mu\) will regulate intermediate
formulas, and \(M\) denotes a soft upper scale below which the eikonal
approximation is valid:
\[
  0<\mu<E_T<M .
\]

\begin{definition}[Resolution-inclusive transition probability]
Let \(S_{\beta;\{k_r,h_r\}|\alpha}\) be the scattering matrix element
with hard charged configuration \(\beta\) and additional photons
\((k_r,h_r)\).  The resolution-inclusive probability at threshold \(E_T\) is
the sum
\[
  P_{\beta|\alpha}(E_T)
  =
  \sum_{N=0}^{\infty}\frac1{N!}
  \sum_{h_1,\ldots,h_N}
  \int_{\mathcal R_{E_T}}
  \prod_{r=1}^{N}
  \frac{\dd^3\vec k_r}{(2\pi)^3\,2\omega_r}\,
  \left|
    S_{\beta;\{k_r,h_r\}|\alpha}
  \right|^2 ,
\]
where \(\mathcal R_{E_T}\) is the detector-degenerate soft region.  At leading
soft logarithmic order one may take \(\mu<\omega_r<E_T\) independently for
each unresolved photon.  Refining the detector model changes finite functions
of \(E_T/M\), but not the cancellation of the regulator \(\mu\).
\end{definition}

\section{External-Line Soft Factors}

In scalar QED, attach an outgoing photon of momentum \(k\) and polarization
vector \(e_h^\mu(k)\) to an outgoing external scalar line of momentum \(p\).
The factor multiplying the hard amplitude is
\[
  (-\ii)\,
  \frac{1}{(p+k)^2+m^2-\ii\epsilon}\,
  \ii g(2p+k)^\mu e^*_{h,\mu}(k).
\]
Since \(p^2+m^2=0\) and \(k^2=0\), the leading term as \(k\to0\) is
\[
  g\,\frac{p\cdot e_h^*(k)}{p\cdot k-\ii\epsilon}.
\]
The denominator is linear in the photon energy.  The soft photon therefore
couples to the charge and momentum of the external line, while the internal
structure of the hard subdiagram is unresolved at leading order.

For the hard transition \(\alpha\to\beta\), the one-soft-photon amplitude is
\[
  \mathcal M_{\beta;k,h|\alpha}
  =
  \mathcal M_{\beta|\alpha}
  \left[
    \sum_n
    \frac{g_n\,p_n\cdot e_h^*(k)}
         {\eta_n p_n\cdot k-\ii\epsilon}
  \right]
  +O(k^0).
\]
The same leading factor is obtained in spinor QED.  The numerator of the
nearly on-shell fermion propagator acts on the external spinor and reduces to
the corresponding momentum factor; spin first enters in subleading soft
terms.
Indeed, for an outgoing fermion line the soft attachment contains
\[
  \left[
  \bar u^\sigma(p)(-g\gamma^\mu)
  \frac{-\ii[-\ii(\not p+\not k)+m]}
       {(p+k)^2+m^2-\ii\epsilon}
  \right]e^*_{h,\mu}(k).
\]
The denominator is \(2p\cdot k-\ii\epsilon+O(k^2)\).  In the numerator,
using the on-shell equation for the external spinor and moving
\(\gamma^\mu\) through the Dirac numerator gives
\[
  \bar u^\sigma(p)\gamma^\mu(-\ii\not p+m)
  =
  \bar u^\sigma(p)(\ii\not p+m)\gamma^\mu
  -2\ii p^\mu\bar u^\sigma(p)
  =
  -2\ii p^\mu\bar u^\sigma(p).
\]
Thus the leading soft limit is
\[
  g\,\frac{p\cdot e_h^*(k)}{p\cdot k-\ii\epsilon}\,
  \bar u^\sigma(p),
\]
the same eikonal factor multiplying the original spinor amplitude.

For \(N\) real soft photons the leading factors multiply:
\[
  \mathcal M_{\beta;\{k_r,h_r\}|\alpha}
  =
  \mathcal M_{\beta|\alpha}
  \prod_{r=1}^{N}
  \left[
    \sum_n
    \frac{g_n\,p_n\cdot e_{h_r}^*(k_r)}
         {\eta_n p_n\cdot k_r-\ii\epsilon}
  \right]
  +O(k_r^0).
\]
For two photons emitted from the same line, the two possible orderings of the
emissions turn the nested eikonal denominators into the product displayed
above.  For an outgoing scalar line this is the elementary identity
\[
\begin{aligned}
  &\frac{g\,p\cdot e_{h_2}^*(k_2)}
        {p\cdot k_2-\ii\epsilon}\,
   \frac{g\,p\cdot e_{h_1}^*(k_1)}
        {p\cdot(k_1+k_2)-\ii\epsilon}
  +(1\leftrightarrow2)
\\
  &\qquad =
  \frac{g\,p\cdot e_{h_1}^*(k_1)}
       {p\cdot k_1-\ii\epsilon}\,
  \frac{g\,p\cdot e_{h_2}^*(k_2)}
       {p\cdot k_2-\ii\epsilon}.
\end{aligned}
\]
Corrections from the non-eikonal parts of the propagators and vertices are
less singular in the soft scaling.  The same combinatorics gives the product
for any finite number of soft photons.

\begin{figure}[!ht]
  \centering
  \resizebox{0.96\textwidth}{!}{%
  \begin{tikzpicture}[>=Stealth,x=1cm,y=1cm,every node/.style={font=\small}]
    \tikzset{
      photon/.style={
        purple!80!black,
        decorate,
        decoration={snake,amplitude=0.75mm,segment length=3.2mm},
        line width=0.95pt,
        -{Latex[length=1.6mm]}
      },
      matter/.style={line width=1.0pt,-{Latex[length=1.8mm]}},
      blob/.style={circle,draw=black,fill=gray!18,line width=0.9pt,minimum size=0.86cm}
    }
    \begin{scope}
      \node[font=\normalsize] at (0,2.3) {one soft photon};
      \node[blob] (B) at (-1.75,0) {};
      \draw[matter] (-3.7,0)--(B.west);
      \draw[matter] (B.east)--(0.35,0);
      \draw[matter] (0.35,0)--(2.4,0);
      \draw[photon] (0.35,0)--(1.18,1.12);
      \node[blue!70!black,below] at (-0.38,-0.12) {\(p+k\)};
      \node[blue!70!black,below] at (1.5,-0.12) {\(p\)};
      \node[purple!80!black,above right] at (1.13,1.05) {\(k,h\)};
      \node[align=center] at (-1.75,-0.92) {hard\\process};
      \node[align=center,blue!65!black] at (0.8,-1.3)
        {\(g\,p\cdot e_h^*/(p\cdot k-\ii\epsilon)\)};
    \end{scope}

    \begin{scope}[shift={(6.2,0)}]
      \node[font=\normalsize] at (0,2.3) {two soft photons on one line};
      \node[blob,minimum size=0.75cm] (C) at (-2.0,0) {};
      \draw[matter] (-3.55,0)--(C.west);
      \draw[matter] (C.east)--(2.5,0);
      \draw[photon] (-0.6,0)--(-0.2,1.05);
      \draw[photon] (0.6,0)--(1.0,1.05);
      \node[purple!80!black] at (-0.12,1.25) {\(k_1,h_1\)};
      \node[purple!80!black] at (1.16,1.25) {\(k_2,h_2\)};
      \node[font=\large] at (2.92,0.06) {\(+\)};
      \begin{scope}[shift={(3.55,0)}]
        \node[blob,minimum size=0.62cm] (D) at (-0.6,0) {};
        \draw[matter] (-1.85,0)--(D.west);
        \draw[matter] (D.east)--(2.0,0);
        \draw[photon] (0.05,0)--(0.42,1.0);
        \draw[photon] (1.08,0)--(1.45,1.0);
      \end{scope}
      \node[align=center,blue!65!black] at (1.5,-1.2)
        {sum over orderings\\equals product of one-photon factors};
    \end{scope}
  \end{tikzpicture}
  }
  \caption{Soft factorization on external charged lines.  A photon with
  wavelength much larger than the hard interaction region resolves only the
  external charge and momentum.}
  \label{fig:volume-ii-soft-factorization}
\end{figure}

\section{Real Soft Emission}

Divide the real-emission contribution by the hard rate with no unresolved
photons.  Using the leading soft factor gives
\[
\begin{aligned}
  R_{\rm real}(E_T,\mu)
  &=
  \sum_{N=0}^{\infty}\frac1{N!}
  \prod_{r=1}^{N}
  \int_{\mu<\omega_r<E_T}
  \frac{\dd^3\vec k_r}{(2\pi)^3\,2\omega_r}
  \sum_{h_r=\pm}
  \left|
    \sum_n
    \frac{g_n\,p_n\cdot e_{h_r}^*(k_r)}
         {\eta_n p_n\cdot k_r}
  \right|^2 .
\end{aligned}
\]
For a single soft photon, define the nonnegative angular kernel
\[
  \mathcal S(k)
  =
  \sum_{h=\pm}
  \left|
    \sum_n
    \frac{g_n\,p_n\cdot e_h^*(k)}
         {\eta_n p_n\cdot k}
  \right|^2 .
\]
A covariant physical-polarization sum may be written as
\[
  \sum_{h=\pm}e_{h,\mu}(k)e^*_{h,\nu}(k)
  =
  \eta_{\mu\nu}+k_\mu c_\nu+k_\nu c_\mu,
  \qquad
  k\cdot c=-1,\qquad c^2=0 .
\]
The terms involving \(c_\mu\) vanish after using
\(\sum_n\eta_n g_n=0\).  Thus
\[
  \mathcal S(k)
  =
  \sum_{n,m}
  \frac{\eta_n\eta_m g_ng_m\,p_n\cdot p_m}
       {(p_n\cdot k)(p_m\cdot k)} .
\]

Writing \(k^\mu=\omega(1,\hat k)\), the energy integral separates from the
angular integral:
\[
  \int_{\mu}^{E_T}
  \frac{\dd^3\vec k}{(2\pi)^3\,2\omega}\,\mathcal S(k)
  =
  \frac{1}{2(2\pi)^3}
  \int_{\mu}^{E_T}\frac{\dd\omega}{\omega}
  \int\dd^2\hat k
  \sum_{n,m}
  \frac{\eta_n\eta_m g_ng_m\,p_n\cdot p_m}
       {(p_n^0-\vec p_n\cdot\hat k)
        (p_m^0-\vec p_m\cdot\hat k)} .
\]
For massive external particles, define the relative-velocity parameter
\[
  \beta_{nm}
  =
  \sqrt{1-\frac{m_n^2m_m^2}{(p_n\cdot p_m)^2}} .
\]
The diagonal term \(n=m\) is understood as the \(\beta_{nm}\to0\) limit.  The
angular integral may be evaluated directly by combining the two denominators:
\[
\begin{aligned}
  &\int\dd^2\hat k\,
  \frac{1}
       {(p_n^0-\vec p_n\cdot\hat k)
        (p_m^0-\vec p_m\cdot\hat k)}
\\
  &=
  \int_0^1\dd x
  \int\dd^2\hat k\,
  \frac{1}
  {\left[
    p_n^0x+p_m^0(1-x)
    -(\vec p_nx+\vec p_m(1-x))\cdot\hat k
  \right]^2}.
\end{aligned}
\]
Choose the polar axis along
\(\vec p_nx+\vec p_m(1-x)\).  The \(\hat k\)-integral becomes an elementary
integral of
\[
  \frac{2\pi\sin\theta\,\dd\theta}
  {\left[
    p_n^0x+p_m^0(1-x)
    -|\vec p_nx+\vec p_m(1-x)|\cos\theta
  \right]^2}.
\]
Performing the \(x\)-integral gives
\[
  \int\dd^2\hat k\,
  \frac{1}
       {(p_n^0-\vec p_n\cdot\hat k)
        (p_m^0-\vec p_m\cdot\hat k)}
  =
  -\frac{2\pi}{p_n\cdot p_m}\,
  \frac{1}{\beta_{nm}}
  \log\frac{1+\beta_{nm}}{1-\beta_{nm}} .
\]
It is useful to introduce
\[
  A_{\beta\alpha}
  =
  -\frac{1}{8\pi^2}
  \sum_{n,m}
  \frac{\eta_n\eta_m g_ng_m}{\beta_{nm}}
  \log\frac{1+\beta_{nm}}{1-\beta_{nm}} .
\]
The representation by \(\mathcal S(k)\) shows \(A_{\beta\alpha}\ge0\) for the
massive charged scattering considered here.  The real unresolved photons
therefore give
\[
  R_{\rm real}(E_T,\mu)
  =
  \sum_{N=0}^{\infty}
  \frac1{N!}
  \left[
    A_{\beta\alpha}\log\frac{E_T}{\mu}
  \right]^N
  =
  \left(\frac{E_T}{\mu}\right)^{A_{\beta\alpha}}
\]
at leading soft logarithmic order.

\section{Virtual Soft Photons and External Factors}

A virtual soft photon exchanged between two distinct external charged lines
also factorizes.  For external lines \(n\) and \(m\), set
\[
  J_{nm}
  =
  -\ii\,p_n\cdot p_m
  \int_{\mu<|\vec k|<M}
  \frac{\dd^4k}{(2\pi)^4}\,
  \frac{1}
       {(k^2-\ii\epsilon)
        (\eta_n p_n\cdot k-\ii\epsilon)
        (-\eta_m p_m\cdot k-\ii\epsilon)} .
\]
The virtual soft photons exchanged among distinct external lines exponentiate:
\[
  \exp\!\left[
    \frac12\sum_{n\ne m}g_ng_mJ_{nm}
  \right].
\]
This expression omits the soft photons whose two endpoints lie on the same
external charged line.  Those terms are tied to the infrared part of the LSZ
normalization factors.  A direct way to see the accounting is to insert the
electromagnetic current
\[
  \widehat j^\mu(x)=\partial_\nu \widehat F^{\mu\nu}(x)
\]
between one-particle states.  In scalar QED, with
\(k=p_2-p_1\) and
\(\omega_p=\sqrt{\vec p^{\,2}+m^2}\),
\[
  \langle \vec p_2|\widehat j^\mu(0)|\vec p_1\rangle
  =
  \frac{g}{(2\pi)^3\sqrt{2\omega_{p_1}}\sqrt{2\omega_{p_2}}}
  (p_1+p_2)^\mu F(k^2),
  \qquad F(0)=1 .
\]
In spinor QED the corresponding decomposition is
\[
\begin{aligned}
  &\langle \vec p_2,\sigma_2|\widehat j^\mu(0)
  |\vec p_1,\sigma_1\rangle
\\
  &\qquad =
  \frac{\ii g}{(2\pi)^3}\,
  \bar u^{\sigma_2}(\vec p_2)
  \left[
    \gamma^\mu F(k^2)
    -\frac{\ii}{2m}(p_1+p_2)^\mu G(k^2)
  \right]
  u^{\sigma_1}(\vec p_1),
  \qquad F(0)+G(0)=1 .
\end{aligned}
\]
In these conventions the one-loop value \(G(0)=-g^2/(8\pi^2)+O(g^4)\) is a
finite contribution to the magnetic form factor.  The electromagnetic matrix
element at \(k^2=0\) is free of infrared divergence, although derivatives such
as \(\left.\dd F/\dd k^2\right|_{k^2=0}\) need not be.  Thus the soft virtual
photons attached to the current insertion must be accompanied by the same
infrared part of the LSZ factors.
In the current form factor, soft photons connecting
the two charged legs exponentiate as \(e^{g^2J_{12}}\).  Finiteness at
\(k^2=0\) then fixes the infrared factor in \(Z\):
\[
  Z_{\rm IR}
  =
  \left.
  e^{-g^2J_{12}}
  \right|_{p_2=p_1,\ \eta_2=-\eta_1}.
\]

Equivalently, compare this with the same-line expression obtained by setting
the two eikonal velocities equal before interpreting the self-contraction:
\[
  ``J_{11}''
  =
  -\ii p_1^2
  \int_{\mu<|\vec q|<M}
  \frac{\dd^4q}{(2\pi)^4}
  \frac{1}
       {(q^2-\ii\epsilon)
        (p_1\cdot q-\ii\epsilon)
        (-p_1\cdot q-\ii\epsilon)} .
\]
The two matter poles in the \(q^0\)-plane contribute to the imaginary part.
The real part satisfies
\[
  \operatorname{Re}J_{11}
  =
  -\left.
  \operatorname{Re}J_{12}
  \right|_{p_2=p_1,\ \eta_2=-\eta_1},
  \qquad
  Z_{\rm IR}=e^{g^2\operatorname{Re}J_{11}} .
\]

At one-loop order this same statement follows directly from the scalar-QED
self-energy.  The soft part is
\[
  \ii\Sigma(p)
  =
  \int\frac{\dd^4k}{(2\pi)^4}
  \frac{-\ii}{(p+k)^2+m^2-\ii\epsilon}
  \frac{-\ii}{k^2-\ii\epsilon}
  (\ii g)^2(2p+k)^2+\cdots ,
\]
whose infrared approximation is
\[
  \ii\Sigma(p)\big|_{\rm soft}
  =
  4g^2p^2
  \int_{|\vec k|<M}
  \frac{\dd^4k}{(2\pi)^4}
  \frac{1}{p^2+m^2+2p\cdot k-\ii\epsilon}
  \frac{1}{k^2-\ii\epsilon}.
\]
The corresponding external factor
\[
  Z
  =
  \left(
    1-\left.
    \frac{\partial\Sigma}{\partial p^2}\right|_{p^2=-m^2}
  \right)^{-1}
\]
contains the same soft logarithm as the same-line eikonal loop.  In the
infrared region,
\[
  \log Z_{\rm IR}
  =
  \ii g^2\,4p^2
  \int_{|\vec k|<M}
  \frac{\dd^4k}{(2\pi)^4}
  \frac{1}
       {(2p\cdot k-\ii\epsilon)^2(k^2-\ii\epsilon)}
  =
  g^2\operatorname{Re}J_{11}+O(g^4).
\]
Thus soft loops beginning and ending on a single external charged line are
not additional uncancelled infrared factors; they are already included in
the LSZ normalization of that charged external line.

\begin{figure}[!ht]
  \centering
  \resizebox{0.82\textwidth}{!}{%
  \begin{tikzpicture}[>=Stealth,x=1cm,y=1cm,every node/.style={font=\small}]
    \tikzset{
      photon/.style={purple!80!black,decorate,
        decoration={snake,amplitude=0.72mm,segment length=3.1mm},
        line width=0.95pt},
      matter/.style={line width=1.0pt,-{Latex[length=1.8mm]}},
      blob/.style={circle,draw=black,fill=gray!18,line width=0.9pt,minimum size=0.75cm},
      zblob/.style={circle,draw=purple!80!black,fill=purple!12,line width=0.9pt,
        minimum size=0.42cm}
    }
    \begin{scope}
      \node[font=\normalsize] at (0,2.05) {current form factor};
      \draw[matter] (-3.3,0)--(-0.45,0);
      \node[blob] at (0,0) {};
      \draw[matter] (0.45,0)--(3.3,0);
      \draw[photon] (0,-1.25)--(0,-0.42);
      \draw[photon] (-2.05,0) .. controls (-1.55,1.05) and
        (1.25,1.05) .. (1.85,0);
      \node[purple!80!black] at (-0.1,1.28) {soft virtual photons};
      \node[blue!70!black] at (-2.65,-0.35) {\(Z^{1/2}\)};
      \node[blue!70!black] at (2.65,-0.35) {\(Z^{1/2}\)};
      \node[purple!80!black,below] at (0,-1.25) {\(j^\mu\)};
    \end{scope}
    \begin{scope}[shift={(7.3,0)}]
      \node[font=\normalsize] at (0,2.05) {same-line soft loops};
      \draw[matter] (-3.25,0)--(3.25,0);
      \node[zblob] at (-1.65,0) {};
      \node[zblob] at (-0.55,0) {};
      \draw[photon] (-0.55,0) .. controls (-0.15,0.85) and
        (1.3,0.85) .. (1.7,0);
      \draw[photon] (0.2,0) .. controls (0.45,0.45) and
        (1.0,0.45) .. (1.25,0);
      \node[purple!80!black] at (-1.1,0.62) {\(Z_{\rm IR}\)};
      \node[purple!80!black] at (1.35,1.1) {absorbed into \(Z_{\rm IR}\)};
      \node[align=center,blue!70!black] at (0,-0.78)
        {no separate same-line\\infrared factor remains};
    \end{scope}
  \end{tikzpicture}
  }
  \caption{The infrared part of the external LSZ factor accounts for soft
  virtual photons whose two endpoints lie on the same charged external line.}
  \label{fig:volume-ii-same-line-ir-cancellation}
\end{figure}

Combining the distinct-line soft loops with the infrared parts of the
\(Z_n^{1/2}\) factors gives
\[
  \exp\!\left[
    \frac12\sum_{n\ne m}g_ng_mJ_{nm}
  \right]\prod_n (Z_n^{\rm IR})^{1/2}
  =
  (\hbox{phase})\,
  \exp\!\left[
    \frac12\sum_{n,m}g_ng_m\,\operatorname{Re}J_{nm}
  \right].
\]

The \(k^0\)-plane poles of \(J_{nm}\), including only the denominators that
control the contour placement, are
\[
  -|\vec k|+\ii\epsilon,\qquad
  |\vec k|-\ii\epsilon,\qquad
  \frac{\vec p_n\cdot\vec k}{p_n^0}-\ii\eta_n\epsilon,\qquad
  \frac{\vec p_m\cdot\vec k}{p_m^0}+\ii\eta_m\epsilon .
\]
The photon poles give the real part of the eikonal integral, while the
matter poles may contribute to its phase.  Evaluating the photon-pole
contribution gives
\[
  \operatorname{Re}J_{nm}
  =
  -\eta_n\eta_m\,
  \frac{p_n\cdot p_m}{2}
  \int_{\mu<|\vec k|<M}
  \frac{\dd^3\vec k}{(2\pi)^3}
  \frac{1}
       {|\vec k|^3
        (p_n^0-\vec p_n\cdot\hat k)
        (p_m^0-\vec p_m\cdot\hat k)} .
\]
Using the same angular integral as for real emission,
\[
  \operatorname{Re}J_{nm}
  =
  \eta_n\eta_m\,
  \frac{1}{8\pi^2\beta_{nm}}
  \log\frac{1+\beta_{nm}}{1-\beta_{nm}}
  \log\frac{M}{\mu}.
\]
Therefore the virtual factor in the amplitude has modulus
\[
  \left(\frac{\mu}{M}\right)^{A_{\beta\alpha}/2},
\]
and its contribution to a rate is
\[
  R_{\rm virt}(M,\mu)
  =
  \left(\frac{\mu}{M}\right)^{A_{\beta\alpha}} .
\]

\begin{figure}[!ht]
  \centering
  \resizebox{0.96\textwidth}{!}{%
  \begin{tikzpicture}[>=Stealth,x=1cm,y=1cm,every node/.style={font=\small}]
    \tikzset{
      photon/.style={purple!80!black,decorate,
        decoration={snake,amplitude=0.72mm,segment length=3.1mm},
        line width=0.95pt},
      matter/.style={line width=1.0pt,-{Latex[length=1.8mm]}},
      blob/.style={circle,draw=black,fill=gray!18,line width=0.9pt,minimum size=0.78cm}
    }
    \begin{scope}
      \node[font=\normalsize] at (0,2.45) {real unresolved photons};
      \node[blob] (A) at (-1.6,0.2) {};
      \draw[line width=1.0pt] (-3.1,0.95)--(A);
      \draw[line width=1.0pt] (-3.1,-0.45)--(A);
      \draw[matter] (A)--(1.35,1.1);
      \draw[matter] (A)--(1.35,-0.68);
      \foreach \x/\y/\lab in {-0.55/0.18/{k_1},0.1/0.02/{k_2},0.63/0.45/{k_N}} {
        \draw[photon,-{Latex[length=1.45mm]}] (\x,\y)--(\x+0.42,\y+0.65);
        \node[purple!80!black] at (\x+0.57,\y+0.75) {\(\lab\)};
      }
      \node[align=center,blue!65!black] at (-0.25,-1.35)
        {\(\displaystyle R_{\rm real}
        =\left(E_T/\mu\right)^{A_{\beta\alpha}}\)};
    \end{scope}

    \begin{scope}[shift={(5.2,0)}]
      \node[font=\normalsize] at (0,2.45) {virtual soft cloud};
      \draw[line width=1.0pt] (-2.35,1.0)--(2.35,1.22);
      \draw[line width=1.0pt] (-2.35,-0.72)--(2.35,-0.94);
      \foreach \x in {-1.25,-0.28,0.66} {
        \draw[photon] (\x,1.03)--(\x+0.34,-0.82);
      }
      \draw[photon] (1.38,1.16) .. controls (2.05,0.45) and (1.93,-0.12) .. (1.55,-0.92);
      \node[align=center] at (0,-1.52)
        {soft loops plus \(Z_n^{1/2}\)};
      \node[align=center,blue!65!black] at (0,-2.28)
        {\(\displaystyle R_{\rm virt}
        =\left(\mu/M\right)^{A_{\beta\alpha}}\)};
    \end{scope}

    \begin{scope}[shift={(10.4,0)}]
      \node[font=\normalsize] at (0,2.45) {inclusive rate};
      \draw[rounded corners=2pt,fill=green!8,draw=green!45!black,line width=0.9pt]
        (-2.15,-0.85) rectangle (2.15,1.15);
      \node[align=center] at (0,0.73)
        {same observed hard event\\for photons below \(E_T\)};
      \draw[photon] (-1.05,-0.55)--(-0.48,-0.05);
      \draw[photon] (0.18,-0.58)--(0.82,-0.12);
      \node[align=center,blue!65!black] at (0,-1.5)
        {\(\displaystyle
        \left(\mu/M\right)^{A_{\beta\alpha}}
        \left(E_T/\mu\right)^{A_{\beta\alpha}}
        =
        \left(E_T/M\right)^{A_{\beta\alpha}}\)};
    \end{scope}
  \end{tikzpicture}
  }
  \caption{Real unresolved photons and virtual soft photons carry opposite
  powers of the auxiliary infrared cutoff.  The product is finite for a fixed
  detector threshold.}
  \label{fig:volume-ii-bloch-nordsieck-cancellation}
\end{figure}

\section{Inclusive Probabilities}

The cancellation of the auxiliary photon-mass regulator between real
unresolved emission and virtual soft exchange is the Bloch--Nordsieck
cancellation.  Equivalently, real soft emission cancels virtual infrared
divergences at the level of inclusive probabilities.  Multiplying the real
and virtual factors gives the leading soft result
\[
  P_{\beta|\alpha}(E_T)
  =
  P_{\beta|\alpha}^{\rm hard}(M)\,
  \left(\frac{E_T}{M}\right)^{A_{\beta\alpha}}
  \times
  F_{\beta\alpha}(E_T/M),
\]
where \(F_{\beta\alpha}\) is finite as \(\mu\to0\) and equals \(1\) in the
leading logarithmic approximation.  The scale \(M\) separates hard and soft
physics; dependence on it cancels once the hard finite correction is computed
to the same order.

For a fixed finite number \(N\) of photons, the rate contains
\[
  \left(\frac{\mu}{M}\right)^{A_{\beta\alpha}}
  \frac1{N!}
  \left[
    A_{\beta\alpha}\log\frac{E_T}{\mu}
  \right]^N ,
\]
which tends to zero as \(\mu\to0\) when \(A_{\beta\alpha}>0\).  The
resolution-inclusive sum remains finite because arbitrarily many unresolved
soft photons are included.

\begin{remark}
The fixed-photon-number matrix elements are exclusive kernels for an idealized
measurement that specifies the exact number of photons down to zero energy.
The inclusive probability above is the transition probability associated with
a finite energy resolution.  Both statements follow from the same soft
factorization formulas.
\end{remark}

The discussion so far treated massive charged particles, where the angular
integral is finite.  If massless charged particles occur, photon emission
collinear to a charged external line produces additional singularities.  The
Kinoshita--Lee--Nauenberg construction enlarges the inclusive sum: one sums
over final states and averages over initial states inside the same detector
degeneracy class.  In symbols, if \([\alpha]\) and \([\beta]\) are the
incoming and outgoing equivalence classes defined by the detector resolution,
one considers
\[
  P([\alpha]\to[\beta])
  =
  \sum_{\alpha'\in[\alpha]}\omega_{\alpha'}
  \sum_{\beta'\in[\beta]}
  \left|S_{\beta'\alpha'}\right|^2,
\]
with nonnegative weights \(\omega_{\alpha'}\) normalized on the incoming
degenerate class.  Soft degeneracy and collinear degeneracy are therefore
handled by the same operational rule: the transition probability is attached
to what the detector distinguishes.

\section{Soft Photon Number Distribution}

The leading soft formula also determines the unresolved photon multiplicity
distribution.  Let
\[
  \overline N_{\beta\alpha}(\mu,E_T)
  =
  A_{\beta\alpha}\log\frac{E_T}{\mu}
\]
be the mean number of photons in the energy window
\(\mu<\omega<E_T\) generated by the eikonal current.  The probability for
emitting exactly \(N\) photons in this unresolved window has the Poisson form
\[
  P_N(\mu,E_T)
  =
  \exp\{-\overline N_{\beta\alpha}\}
  \frac{\overline N_{\beta\alpha}^{\,N}}{N!}
\]
at leading soft order, multiplied by the hard probability and finite
corrections.  The exponential factor is the virtual soft factor together with
the no-resolved-emission probability.  Summing over \(N\) gives unity for the
soft probability distribution, while the hard probability is multiplied by the
finite detector-resolution factor
\[
  \left(\frac{E_T}{M}\right)^{A_{\beta\alpha}} .
\]

As \(\mu\to0\), the mean \(\overline N_{\beta\alpha}\) diverges when
\(A_{\beta\alpha}>0\).  The inclusive probability is finite because the
detector state is a distribution over arbitrarily many unresolved photons.
The exclusive fixed-\(N\) probability tends to zero because it samples a set
of vanishing measure inside the infrared cloud selected by the long-range
electromagnetic field.

\section{Dressed Charged States}

The same eikonal data can also be incorporated into the asymptotic states.
For a charged particle of momentum \(p\), a soft coherent dressing has the
leading soft form
\[
  \mathcal W_p
  =
  \exp\!\left[
    \sum_h
    \int_{\omega<\Lambda}
    \frac{\dd^3\vec k}{(2\pi)^3\,2\omega}
    \left(
      f_p^h(k)a_h^\dagger(k)
      -
      \overline{f_p^h(k)}a_h(k)
    \right)
  \right],
\]
with leading coefficient
\[
  f_p^h(k)\sim g\,\frac{p\cdot e_h^*(k)}{p\cdot k}.
\]
The operator \(\mathcal W_p\) attaches the long-range electromagnetic field
carried by the charged particle.  This is the leading Faddeev--Kulish
dressing datum: the asymptotic charged state is a charged hard state together
with its coherent soft electromagnetic field.  Dressed-state and inclusive
descriptions use the same soft factor.  The inclusive construction sums over
unresolved radiation in probabilities, while the dressed construction builds
part of the soft cloud into the state space.

\section{A Gauge-Invariant Infrared Regulator}

The cutoff \(\mu\) in the formulas above can be realized by deforming the
theory so that the photon has a small mass, and then removing the deformation.
If the same soft region is regulated dimensionally instead, the logarithms
\(\log(\mu/M)\) are replaced by infrared poles \(1/\epsilon_{\rm IR}\) plus
scheme-dependent finite constants; the coefficient multiplying the logarithm
or pole is the same eikonal coefficient \(A_{\beta\alpha}\), so the
real--virtual cancellation is identical coefficient by coefficient.
A convenient gauge-invariant deformation is an Abelian-Higgs regulator.
Introduce a complex scalar \(\varphi\), charged under the gauge field
\(A_\mu\), with
\[
  D_\mu\varphi=(\partial_\mu+\ii g_H A_\mu)\varphi
\]
and Lagrangian terms
\[
  \mathcal L_{\rm reg}
  =
  -\frac14F_{\mu\nu}F^{\mu\nu}
  -|D_\mu\varphi|^2
  -V(|\varphi|).
\]
Choose the potential so that its vacuum satisfies
\[
  |\vev{\varphi}|=v\ne0 .
\]
Expanding around this vacuum gives
\[
  -|D_\mu\varphi|^2
  \supset
  -g_H^2v^2 A_\mu A^\mu ,
\]
so the gauge boson acquires a mass proportional to \(g_Hv\).  Taking \(v\) to
zero after forming inclusive quantities gives a consistent version of the
infrared-regulated calculation.

\begin{figure}[!ht]
  \centering
  \resizebox{0.9\textwidth}{!}{%
  \begin{tikzpicture}[>=Stealth,x=1cm,y=1cm,every node/.style={font=\small}]
    \begin{scope}
      \node[font=\normalsize] at (0,2.45) {regulator scalar};
      \draw[->,line width=0.85pt] (-3.0,0)--(3.1,0)
        node[below] {\(\operatorname{Re}\varphi\)};
      \draw[->,line width=0.85pt] (0,-0.28)--(0,2.1)
        node[left] {\(V\)};
      \draw[blue!70!black,very thick,domain=-2.35:2.35,samples=120]
        plot (\x,{0.23*(\x*\x-1.55)^2+0.08});
      \fill[green!50!black] (1.245,0.08) circle (2.4pt);
      \node[green!45!black,below] at (1.245,-0.12) {\(v\)};
      \node[align=center,blue!70!black] at (0,-0.75)
        {choose a vacuum\\\(|\vev{\varphi}|=v\)};
    \end{scope}

    \begin{scope}[shift={(6.3,0)}]
      \node[font=\normalsize] at (0,2.45) {massive photon regulator};
      \draw[line width=1.0pt] (-2.3,0.72)--(-0.6,0.72);
      \draw[purple!80!black,decorate,
        decoration={snake,amplitude=0.78mm,segment length=3.2mm},
        line width=0.95pt]
        (-0.6,0.72)--(1.25,0.72);
      \draw[line width=1.0pt] (1.25,0.72)--(2.4,0.72);
      \node[purple!80!black,below] at (0.34,0.47) {\(m_\gamma\propto g_Hv\)};
      \draw[->,line width=1.0pt] (0,0.3)--(0,-0.62);
      \node[align=center] at (0,-1.15)
        {\(-g_H^2v^2A_\mu A^\mu\)\\from \(|D_\mu\varphi|^2\)};
      \node[align=center,orange!80!black] at (0,1.78)
        {compute at \(v\ne0\),\\then take \(v\to0\)};
    \end{scope}
  \end{tikzpicture}
  }
  \caption{A small photon mass may be introduced through a Higgs deformation
  while keeping gauge invariance manifest.}
  \label{fig:volume-ii-abelian-higgs-ir-regulator}
\end{figure}

The infrared analysis will be used later whenever massless gauge fields occur.
The next step in the logical development is the generating functional
formalism, which packages connected Green functions and one-particle
irreducible vertices in a form suitable for renormalization.
